% Reviewed translation: PDF pages 350(part)--355 / printed pages 336(part)--341.
\MathBookSourcePage{350}
\section{上风格式的数值分析}
\label{sec:navier-upwind-analysis}

当黏度 $\nu$ 很小时, 或者等价地, 当 Reynolds 数 $Re$ 与流体的其他参数相比很大时,
在 $\Gamma$ 邻域会出现一个边界层; 黏性在该处占主导, 而在 $\Omega$ 内部则可忽略.
与此同时, 流动变为湍流. 因而, 在高 Reynolds 数下, Navier--Stokes 方程的解具有严重的不连续性.

这里的目的不是对湍流进行建模, 但仍值得考察 Navier--Stokes 方程不连续解的一种可能离散化.
我们不采用迄今研究的中心格式, 而提出更适于描述不连续流动的上风格式. 后面的分析表明,
这些格式几乎是最优的. 读者也可参见 Johnson \& Saranen~\cite{JohnsonSaranen47}
关于另一种方法, 即适用于 Euler 和 Navier--Stokes 方程的流线扩散法.

\subsection{流函数--涡量格式中的上风化}
\label{subsec:navier-upwind-stream-vorticity}

本节讨论的上风格式由 Fortin~\cite{Fortin31} 发展, 其灵感来自
Lesaint \& Raviart~\cite{LesaintRaviart51} 为求解中子输运方程而提出的一种数值方法.
它源自下面的启发式说明. 当光滑向量场 $\boldsymbol u$ 无散时, 对流项满足
(采用重复指标表示求和的通常求和约定)
\[
u_j(\partial u_i/\partial x_j)=\partial(u_ju_i)/\partial x_j.
\]

现在设 $\Omega$ 是两个有界区域 $\Omega_1$ 和 $\Omega_2$ 的并, 它们由
Figure~\ref{fig:navier-upwind-interface} 所示的界面 $S$ 分隔. 以 $\Gamma$ 表示
$\Omega$ 的边界, 以 $\boldsymbol n_i$ 表示 $\Omega_i$ 在 $S$ 上的单位外法向量.

\setcounter{figure}{20}
\begin{figure}[htbp]
  \centering
  \includegraphics[width=0.70\textwidth]{chapter-04-figure-21.png}
  \caption{两个子区域之间的上风界面.}
  \label{fig:navier-upwind-interface}
\end{figure}

\MathBookSourcePage{351}
假设向量场 $\boldsymbol u$ 在 $\Omega$ 中不再整体光滑, 而是满足
\[
\boldsymbol u|_{\Omega_k}\in H^1(\Omega_k)^N,
\qquad \boldsymbol u|_\Gamma=\boldsymbol0,
\qquad \boldsymbol u\in H(\operatorname{div};\Omega),
\qquad \operatorname{div}\boldsymbol u=0.
\]
于是乘积 $u_ju_i$ 在 $\Omega$ 中不再可微, 但它的分布导数具有如下形式:
\[
[\partial(u_ju_i)/\partial x_j]
=\left.\partial(u_ju_i)/\partial x_j\right|_{\Omega_1\cup\Omega_2}
\]
再加上与场 $\boldsymbol u$ 跨界面 $S$ 的跳跃相对应的面分布, 即
\[
\boldsymbol u\cdot\boldsymbol n
\bigl(u_i^{\mathrm{ext}}-u_i^{\mathrm{int}}\bigr)\,ds,
\]
其中 $\boldsymbol n$ 以及外部和内部的概念都相对于同一区域而言. 上风化依赖如下原则:
\[
\textit{$\Omega_k$ 中的流动仅依赖于通过 $S$ 进入的通量.}
\]
根据这一原则, 引入记号
\begin{equation}\label{eq:navier-upwind-inflow-interface}
\partial_-\Omega_k
=\{x\in S;\ \boldsymbol u\cdot\boldsymbol n_k(x)<0\},
\end{equation}
表示流体进入 $\Omega_k$ 的界面 $S$ 的部分, 并把对流项定义为
\begin{equation}\label{eq:navier-upwind-two-region-convection}
\begin{aligned}
\int_\Omega u_j(\partial u_i/\partial x_j)v_i\,dx
\mathrel{\mathop{=}^{\mathrm{def}}}
&\sum_{k=1}^2\int_{\Omega_k}u_j(\partial u_i/\partial x_j)v_i\,dx\\
&+\sum_{k=1}^2\int_{\partial_-\Omega_k}
\boldsymbol u\cdot\boldsymbol n_k
(u_i^{\mathrm{ext}}-u_i^{\mathrm{int}})v_i^{\mathrm{int}}\,ds.
\end{aligned}
\end{equation}
这里 $\boldsymbol u\in H(\operatorname{div};\Omega)$,
$\operatorname{div}\boldsymbol u=0$,
$\boldsymbol u|_{\Omega_k}\in H^1(\Omega_k)^N$, 且
$\boldsymbol v|_{\Omega_k}\in H^1(\Omega_k)^N$.
显然, 当 $\boldsymbol u$ 和 $\boldsymbol v$ 属于 $H^1(\Omega)^N$ 时, 面项消失,
从而恢复熟悉的对流表达式. 但是, 当 $\boldsymbol u$ 和 $\boldsymbol v$ 并非整体属于
$H^1$ 时, 必须把 \eqref{eq:navier-upwind-two-region-convection} 理解为一个定义.

\MathBookSourcePage{352}
我们拟把上述考虑推广到 $\Omega$ 已三角剖分的情形, 其中每个三角形都被视为
$\Omega$ 的一个子区域. 但在此之前, 先回顾 Section~\ref{subsec:navier-centered-stream-vorticity}
中发展的 Navier--Stokes 系统流函数--涡量表述及其逼近框架.

设 $\Omega$ 是 $\mathbb R^2$ 中具有多边形边界 $\Gamma$ 的有界区域, 因而可被完全三角剖分;
并假设 $\Omega$ 是凸的, 从而正则性条件
\eqref{eq:navier-centered-stream-dual-regularity} 成立. 为简化讨论, 暂时略去压力,
只考虑流函数和涡量. 令
\[
Y=L^r(\Omega)^2,
\qquad
X=\{\operatorname{curl}\phi;\ \phi\in\Phi_s\}\times L^2(\Omega),
\]
其中 $s\geq4$ 且 $1/r+1/s=1$. Stokes 算子 $T\in\mathcal L(Y;X)$ 定义为
$T\boldsymbol f=(\operatorname{curl}\psi,\omega)$, 这里
\begin{equation}\label{eq:navier-upwind-stokes-operator}
\left\{
\begin{aligned}
(\operatorname{curl}\omega,\operatorname{curl}\phi)
&=(\boldsymbol f,\operatorname{curl}\phi)
&&\forall\phi\in\Phi_s,\\
(\operatorname{curl}\psi,\operatorname{curl}\mu)
&=(\omega,\mu)
&&\forall\mu\in W^{1,r}(\Omega),
\end{aligned}
\right.
\end{equation}
其中 $\psi$ 和 $\omega$ 通过
\begin{equation}\label{eq:navier-upwind-velocity-vorticity}
\boldsymbol u=\operatorname{curl}\psi,
\qquad
\omega=\operatorname{curl}\boldsymbol u
\end{equation}
与速度 $\boldsymbol u$ 相联系. 像往常一样, 引入 $X$ 的子空间
\begin{equation}\label{eq:navier-upwind-continuous-subspace}
V=\{(\operatorname{curl}\phi,\theta)\in X;
\ \phi\in H_0^2(\Omega),\ \theta=-\Delta\phi\}.
\end{equation}
由于 $\Omega$ 是凸的, $\Gamma$ 只有一个连通分支; 并且由
Theorem~\ref{thm:incompressible-stream-vorticity-wellposed} 已知 $T\in\mathcal L(Y;V)$.
非线性项的表达方式与 \eqref{eq:navier-centered-stream-nonlinearity} 大致相同:
\begin{equation}\label{eq:navier-upwind-nonlinearity}
G(u)=\omega\operatorname{grad}\psi-\boldsymbol f,
\qquad
u=(\operatorname{curl}\psi,\omega)\in X.
\end{equation}
于是 Navier--Stokes 问题如 \eqref{eq:navier-centered-stream-continuous} 那样表述:
给定 $\boldsymbol f\in L^r(\Omega)^2$, 求
$u=(\operatorname{curl}\psi,\omega)\in X$, 使
\begin{equation}\label{eq:navier-upwind-continuous-problem}
\left\{
\begin{aligned}
\nu(\operatorname{curl}\omega,\operatorname{curl}\phi)
+(\omega\operatorname{grad}\psi,\operatorname{curl}\phi)
&=(\boldsymbol f,\operatorname{curl}\phi)
&&\forall\phi\in\Phi_s,\\
(\operatorname{curl}\psi,\operatorname{curl}\mu)
&=(\omega,\mu)
&&\forall\mu\in W^{1,r}(\Omega),
\end{aligned}
\right.
\end{equation}
其中 $\psi$ 和 $\omega$ 由 \eqref{eq:navier-upwind-velocity-vorticity} 与 $\boldsymbol u$
相联系. 取上述 $T,G$ 和 $\lambda=1/\nu$, 可紧凑地写为
\begin{equation}\label{eq:navier-upwind-continuous-operator}
F(\lambda,u)=u+\lambda TG(u)=0.
\end{equation}

为进行逼近, 首先引入由 \eqref{eq:navier-centered-stream-discrete-spaces} 定义的
两个有限维空间 $\Theta_h$ 和 $\Phi_h$:
\[
\Theta_h=\{\theta_h\in\mathcal C^0(\bar\Omega);
\ \theta_h|_\kappa\in P_l\ \forall\kappa\in\mathcal T_h\},
\qquad
\Phi_h=\Theta_h\cap H_0^1(\Omega),
\]
其中 $\mathcal T_h$ 是 $\bar\Omega$ 的一个三角剖分, 且整数 $l\geq1$. 自然取
\[
X_h=\{\operatorname{curl}\phi_h;\ \phi_h\in\Phi_h\}\times\Theta_h\subset X.
\]

\MathBookSourcePage{353}
并定义相应子空间
\begin{equation}\label{eq:navier-upwind-discrete-subspace}
V_h=\{(\operatorname{curl}\phi_h,\theta_h)\in X_h;
(\operatorname{curl}\phi_h,\operatorname{curl}\mu_h)
=(\theta_h,\mu_h)\ \forall\mu_h\in\Theta_h\}.
\end{equation}
离散 Stokes 算子可以定义在比 $Y$ 更宽的空间上, 即
\[
Y_h=\{\operatorname{curl}\phi_h;\ \phi_h\in\Phi_h\}'
\]
配备范数
\[
\|l\|_*=\sup_{\phi_h\in\Phi_h}
\frac{\langle l,\operatorname{curl}\phi_h\rangle}{|\phi_h|_{1,s,\Omega}}.
\]
规定
\[
\langle l,v_h\rangle=\langle l,\operatorname{curl}\phi_h\rangle
\quad\forall v_h=(\operatorname{curl}\phi_h,\theta_h)\in V_h,
\]
还可把空间 $Y_h$ 等同于 $V_h'$ 的一个子空间, 其中 $V_h'$ 是 $V_h$ 的对偶空间;
因此有连续嵌入
\[
Y\hookrightarrow Y_h\hookrightarrow V_h'.
\]
以后还需要为 $Y_h$ 配备 $V_h'$ 的范数, 即令
\[
\|l\|_h=\sup_{v_h\in V_h}\frac{\langle l,v_h\rangle}{\|v_h\|_X}
\quad\forall l\in Y_h;
\]
显然
\[
\|l\|_h\leq\|l\|_*
\quad\forall l\in Y_h.
\]
然后定义逼近 Stokes 算子 $T_h\in\mathcal L(Y_h;X_h)$: 对 $l\in Y_h$,
求 $T_hl=(\operatorname{curl}\psi_h,\omega_h)\in X_h$, 满足
\begin{equation}\label{eq:navier-upwind-discrete-stokes}
\left\{
\begin{aligned}
(\operatorname{curl}\omega_h,\operatorname{curl}\phi_h)
&=\langle l,\operatorname{curl}\phi_h\rangle
&&\forall\phi_h\in\Phi_h,\\
(\operatorname{curl}\psi_h,\operatorname{curl}\mu_h)
&=(\omega_h,\mu_h)
&&\forall\mu_h\in\Theta_h.
\end{aligned}
\right.
\end{equation}
显然, 算子 $T_h$ 的值域就是空间 $V_h$.

现在转向对流项的上风离散化. 把 $\Omega$ 看成 $\mathcal T_h$ 中所有三角形
$\kappa$ 的并, 公式 \eqref{eq:navier-upwind-two-region-convection} 可推广为如下定义:
\begin{equation}\label{eq:navier-upwind-mesh-convection}
\begin{aligned}
\int_\Omega u_j(\partial u_i/\partial x_j)v_i\,dx
\mathrel{\mathop{=}^{\mathrm{def}}}
\sum_{\kappa\in\mathcal T_h}\Bigg\{
&\int_\kappa u_j(\partial u_i/\partial x_j)v_i\,dx\\
&+\int_{\partial_-\kappa}
\boldsymbol u\cdot\boldsymbol n
(u_i^{\mathrm{ext}}-u_i^{\mathrm{int}})v_i^{\mathrm{int}}\,ds
\Bigg\}.
\end{aligned}
\end{equation}
这促使我们拆分对流项, 用两个形式代替单个三线性形式: 一个对应体积分, 一个对应线积分.
因此, 对所有 $u=(\operatorname{curl}\psi,\omega)$ 和
$v=(\operatorname{curl}\phi,\theta)\in X_h$, 以及所有形如
$w_i=\partial\chi/\partial x_j$ 的 $\boldsymbol w$, 其中 $\chi\in\Phi_h$, 定义如下.

\MathBookSourcePage{354}
\begin{equation}\label{eq:navier-upwind-volume-form}
a_1(u;v,\boldsymbol w)
=\sum_{\kappa\in\mathcal T_h}\int_\kappa
u_j(\partial v_i/\partial x_j)w_i\,dx,
\end{equation}
其中 $\boldsymbol u=\operatorname{curl}\psi$, $\boldsymbol v=\operatorname{curl}\phi$;
类似地, 对 $z=(\operatorname{curl}\chi,\mu)\in X_h$, 取
\begin{equation}\label{eq:navier-upwind-boundary-form}
a_2^z(u;v,\boldsymbol w)
=\sum_{\kappa\in\mathcal T_h}\int_{\partial_-\kappa(z)}
\boldsymbol u\cdot\boldsymbol n
(v_i^{\mathrm{ext}}-v_i^{\mathrm{int}})w_i^{\mathrm{int}}\,ds,
\end{equation}
其中 $\boldsymbol n$ 始终表示 $\kappa$ 的单位外法向量, 上标 $\mathrm{ext}$
(相应地 $\mathrm{int}$)表示函数在 $\kappa$ 边界上的外部迹(相应地内部迹), 且
\begin{equation}\label{eq:navier-upwind-element-inflow}
\partial_-\kappa(z)
=\{x\in\partial\kappa;\ (\boldsymbol z\cdot\boldsymbol n)(x)<0\},
\qquad \boldsymbol z=\operatorname{curl}\chi.
\end{equation}
即 $\partial_-\kappa(z)$ 是速度为 $\boldsymbol z$ 的流体进入 $\kappa$ 的边界部分.
最后, 引入形式
\begin{equation}\label{eq:navier-upwind-combined-form}
\widetilde a^z(u;v,\boldsymbol w)
=a_1(u;v,\boldsymbol w)+a_2^z(u;v,\boldsymbol w),
\end{equation}
并由下式定义映射 $G_h:v\in X_h\mapsto G_h(v)\in Y_h$:
\begin{equation}\label{eq:navier-upwind-discrete-nonlinearity}
\langle G_h(v),\operatorname{curl}\chi\rangle
=\widetilde a^v(v;v,\operatorname{curl}\chi)
-(\boldsymbol f,\operatorname{curl}\chi)
\quad\forall\chi\in\Phi_h.
\end{equation}
稍后将建立 $G_h$ 的连续性, 但立即可见, 由于
$a_2^v(v;v,\boldsymbol w)$ 对 $\partial_-\kappa(v)$ 的依赖, 映射 $G_h$ 不可微.
注意, 当 $\boldsymbol w\in H^1(\Omega)^2$ 时, 形式 $\widetilde a$ 满足说明
定义 \eqref{eq:navier-upwind-mesh-convection} 合理性的关键性质:
\begin{equation}\label{eq:navier-upwind-integration-by-parts}
\widetilde a^z(u;v,\boldsymbol w)
=-\int_\Omega u_j(\partial w_i/\partial x_j)v_i\,dx,
\end{equation}
无论函数 $z$ 为何, 当然这里要求
$\operatorname{div}\boldsymbol u=0$ 于 $\Omega$ 且
$\boldsymbol u\cdot\boldsymbol n=0$ 于 $\Gamma$.

于是 Navier--Stokes Problem~\eqref{eq:navier-upwind-continuous-problem}
具有如下上风离散化: 求 $\psi_h\in\Phi_h$ 和 $\omega_h\in\Theta_h$, 满足
\begin{equation}\label{eq:navier-upwind-discrete-problem}
\left\{
\begin{aligned}
\nu(\operatorname{curl}\omega_h,\operatorname{curl}\phi_h)
+\widetilde a^{u_h}(u_h;u_h,\operatorname{curl}\phi_h)
&=(\boldsymbol f,\operatorname{curl}\phi_h)
&&\forall\phi_h\in\Phi_h,\\
(\operatorname{curl}\psi_h,\operatorname{curl}\mu_h)
&=(\omega_h,\mu_h)
&&\forall\mu_h\in\Theta_h,
\end{aligned}
\right.
\end{equation}
其中 $\widetilde a$ 由 \eqref{eq:navier-upwind-volume-form}--
\eqref{eq:navier-upwind-combined-form} 定义. 采用上述记号, 此问题具有
\eqref{eq:navier-branch-nondifferentiable-discrete-map} 的形式:
\begin{equation}\label{eq:navier-upwind-discrete-operator}
F_h(\lambda,u_h)=u_h+\lambda T_hG_h(u_h)=0,
\end{equation}
其中 $u_h=(\operatorname{curl}\psi_h,\omega_h)$ 且 $\lambda=1/\nu$.
注意 $F_h$ 把 $\Lambda\times X_h$ 映到 $V_h$ 中.

\subsection{上风格式的误差分析}
\label{subsec:navier-upwind-error-analysis}

在应用 Subsection~\ref{subsec:navier-nondifferentiable-branches} 的材料建立
\eqref{eq:navier-upwind-discrete-problem} 的收敛性之前, 先证明形式
$\widetilde a(\mathord\cdot;\mathord\cdot,\mathord\cdot)$ 的一些重要性质.
读者会发现, 操作这个形式并不像初看起来那样困难. 首先, 从
$\widetilde a(\mathord\cdot;\mathord\cdot,\mathord\cdot)$ 的表达式可清楚看出,
它要求比空间 $\Phi$ 的 $W^{1,s}$ 范数更精细的范数. 具体而言, 希望采用半范数

\MathBookSourcePage{355}
\[
\left(\sum_{\kappa\in\mathcal T_h}|\phi|_{2,\kappa}^2\right)^{1/2}.
\]
这正是空间 $V_h$ 的用途: Chapter~\ref{chap:incompressible-mixed-stokes} 已经证明,
在适当条件下, $V_h$ 中的函数满足一致稳定性性质
\begin{equation}\label{eq:navier-upwind-uniform-w1t}
|\phi_h|_{1,t,\Omega}\leq C_t\|\theta_h\|_{0,\Omega}
\quad\text{(参见 \eqref{eq:incompressible-sfvp-h1-bound})}
\quad\forall v_h=(\operatorname{curl}\phi_h,\theta_h)\in V_h,
\end{equation}
以及
\begin{equation}\label{eq:navier-upwind-uniform-broken-h2}
\left(\sum_{\kappa\in\mathcal T_h}|\phi_h|_{2,\kappa}^2\right)^{1/2}
\leq C\|\theta_h\|_{0,\Omega}
\quad\text{(参见 \eqref{eq:incompressible-sfvp-auxiliary-h2-bound})},
\end{equation}
其中常数与 $h$ 无关. 换言之:
\[
\textit{$X$ 的半范数 $|v_h|=\|\theta_h\|_{0,\Omega}$ 是 $V_h$ 上的范数,}
\]
并且对于每个 $t\geq2$, 它与
\[
|\phi_h|_{1,t,\Omega}+\|\theta_h\|_{0,\Omega}
\quad\text{以及}\quad
\left(\sum_{\kappa\in\mathcal T_h}|\phi_h|_{2,\kappa}^2\right)^{1/2}
+\|\theta_h\|_{0,\Omega}
\]
都一致等价. 这提示采用如下网格相关范数:
\begin{equation}\label{eq:navier-upwind-mesh-dependent-norm}
[\phi]_t
=\left(\sum_{\kappa\in\mathcal T_h}|\phi|_{2,\kappa}^2\right)^{1/2}
+h^{-2/t}|\phi|_{1,t,\Omega}
\quad\forall t\geq2.
\end{equation}

其次, 由于 $X$ 中的函数正则性很低, 有时采用下面的光滑化算子会很有用:
\[
P\in\mathcal L\bigl(X;
\{\operatorname{curl}\phi;\ \phi\in H^2(\Omega)\cap H_0^1(\Omega)\}
\times L^2(\Omega)\bigr),
\]
定义如下: 对 $u=(\operatorname{curl}\psi,\omega)\in X$,
$Pu=(\operatorname{curl}\phi,\omega)$, 其中 $\phi$ 是 Dirichlet 问题
\begin{equation}\label{eq:navier-upwind-smoothing-dirichlet}
-\Delta\phi=\omega\quad\text{in }\Omega,
\qquad
\phi=0\quad\text{on }\Gamma
\end{equation}
的解. 当 $\Omega$ 为凸域时, $\phi$ 确实属于 $H^2(\Omega)$, 且
\begin{equation}\label{eq:navier-upwind-smoothing-regularity}
\|\phi\|_{2,\Omega}\leq C\|\omega\|_{0,\Omega}.
\end{equation}
显然, 算子 $P$ 在 $V$ 上化为恒等映射. 此外, 当
$u_h=(\operatorname{curl}\psi_h,\omega_h)\in V_h$ 时, 相应函数 $\phi$
与 Lemma~\ref{lem:incompressible-sfvp-hypothesis-h1} 的证明中引入的函数
$\phi(h)$ 相同. 换言之, $u_h=(\operatorname{curl}\psi_h,\omega_h)\in V_h$
和 $Pu_h=(\operatorname{curl}\phi,\omega_h)$ 由
\[
\psi_h=\mathring P_h\phi
\]
联系, 其中 $\mathring P_h$ 是由 \eqref{eq:appendix-dirichlet-ritz-projection} 定义的
$\Phi_h$ 上的 $H^1$ 投影:
\begin{equation}\label{eq:navier-upwind-h1-projection}
(\operatorname{curl}(\mathring P_h\phi-\phi),\operatorname{curl}\mu_h)=0
\quad\forall\mu_h\in\Phi_h.
\end{equation}
下一个 Lemma 给出 $[\phi-\mathring P_h\phi]_t$ 的一个有用估计.
